\section{引言}
\label{sec:introduction}

拒答行为的可干预性为理解语言模型内部表示提供了一个具体研究问题：当部分权重沿某种表示方向改变时，哪些输出行为随之变化，哪些性质仍保持稳定。Arditi 等在所研究的模型与设置中发现，拒答行为可以由一个方向介导，并展示了相应干预\cite{arditi2024}。这一结果支持对拒答相关表示开展机制分析，但不能推导出所有模型、所有语言或所有请求都共享同一条充分的拒答方向。对一个新的模型实例，干预空间与评价协议仍需分别检验。

Heretic 将方向消融组织为自动化搜索流程：从两组提示的残差统计中构造方向，再搜索层位置、组件权重和干预范围，以拒答相关指标与输出分布漂移共同评价候选\cite{heretic}。这种参数化容易解释，也使候选之间能够保持统一接口。其限制在于，局部更新受到所选方向及权重核约束；即便搜索充分，也只覆盖由该表示假设定义的更新族。因此，扩大更新空间是合理的方法动机，但“更大的空间”本身并不意味着更好的实测结果。

任意秩消融（Arbitrary-Rank Ablation，ARA）在上游实现中转向模块输入输出，利用保持、拉近与推远目标求解映射变化\cite{upstreamara}。本文沿用这一思想及名称，把研究对象限定为本地校准式 ARA（Calibrated Arbitrary-Rank Ablation，下文简称本地 ARA；配置字段仍称 CARA）。本地实现通过有限秩因子表达增量，并对距离尺度、推远目标和部署因子进行处理。这里的贡献是对既有思想的适配、明确化和证据分析，不把 ARA 名称、三项基本目标或 LoRA 路径归为本项目首创。

本文区分已经运行的单位置版本 point-v1 与当前实现的轨迹版本 trajectory-v2。前者只采集提示末端预测首 token 的模块状态，后者使用冻结基座短续写的逐步状态，并增加部署增益、连续评分及导出约束。两者不仅采集位置不同，校准规模、验证切分、关键词集合和搜索协议也有差异。因此，当前源码只能说明后续机制已经实现，不能用来重新命名历史结果，更不能把多个协议下的数字汇入同一条效果结论。

实验结论是一个明确的未达标结果。唯一量化来源为 \code{docs/logs/ara-v1/} 的归档：一次搜索包含 120 个候选，全部完成，但无候选通过验收。按验证集关键词命中率优先排序，trial 66 的结果为 54/100，对应首 token KL 为 0.089975；机器验收没有选中该候选。最低 KL 来自另一个候选，二者不能拼接。未修改基座可作为同次运行对照，Heretic 方向消融则没有同协议实测数据。这样的证据支持分析本次搜索的失败位置，不支持判断本地 ARA 优于 Heretic。

本文的工作包含三个相互衔接的部分。方法部分把本地单位置目标写为与源码一致的数学形式，解释有限秩边界、非凸优化和规范化条件。适配部分说明轨迹采集、模块识别与验收状态之间的关系，明确哪些变化尚未获得效果证据。实验部分从原始 journal 重建候选分布，区分进程完成、数值门槛和模型验收，并讨论代理指标与选择偏差。论文的价值在于使方法描述和实验结论具有同一证据边界，而不是通过放宽验收条件把负结果改写为成功。
